%------------------------------------------------------------------------------%
\documentclass[fleqn,utf8,aspectratio=169,12pt]{beamer}

\usepackage{gaal}

\title{Matrizes e Vetores}

%------------------------------------------------------------------------------%
\begin{document}

\addtitle

%------------------------------------------------------------------------------%
\begin{frame}{Matriz}
  \begin{itemize}[<+->]
    \item Uma tabela de números
    \item Definimos operações com matrizes
  \end{itemize}
\end{frame}

%------------------------------------------------------------------------------%
\begin{frame}{Notação}
\[
  A_{3 \times 5}
  \pause
  \;=\;
  \begin{bmatrix}
    80 & 38 & 24 & 8  & 6  \\
    70 & 38 & 20 & 10 & 8  \\
    69 & 38 & 21 & 6  & 11
  \end{bmatrix}
  \pause
  \;=\;
  \begin{pmatrix}
    80 & 38 & 24 & 8  & 6  \\
    70 & 38 & 20 & 10 & 8  \\
    69 & 38 & 21 & 6  & 11
  \end{pmatrix}
\]

\pause
$A$ é o nome da matriz

\pause
3 e 5 são os índices e indicam a ordem da matriz
\end{frame}

%------------------------------------------------------------------------------%array
\begin{frame}{Notação}
Elementos dispostos na horizontal compõem uma \emph{linha} da matriz

\pause
A primeira linha da matriz $A$ é
\(
  [80 \; 38 \; 24 \; 8 \; 6]
\)

\pause
\bigskip
\bigskip
Elementos dispostos na vertical compõem uma \emph{coluna} da matriz

\pause
A quarta coluna da matriz $A$ é
\(
\begin{bmatrix}
  8  \\
  10 \\
  6  \\
  12
\end{bmatrix}
\)
\end{frame}

%------------------------------------------------------------------------------%
\begin{frame}{Notação}
  Elementos da matriz são denotados por $a_{ij}$

  \pause
  \[
    a_{13} = 24\qquad
    a_{31} = 69\qquad
    a_{44} = 12
  \]

  \pause
  \bigskip
  A matriz pode ser denotada por
  \[
     A = [a_{ij}]_{m \times n}
  \]
\end{frame}

%------------------------------------------------------------------------------%
\begin{frame}{Matrizes Especiais}
\begin{description}[<+->]
  \item[Matriz Quadrada]   \tabto{13mm} matriz com o mesmo número de linhas e colunas
                                       \onslide<+->{\emph{$A_n = A_{n \times n}$}}
  \item[Matrix Coluna]     \tabto{13mm} matriz com uma única coluna $B_{1\times n}$
  \item[Matriz Linha]      \tabto{13mm} matriz com uma única linha $C_{m\times 1}$
  \item[Matriz Nula]       \tabto{13mm} matriz \emph{$\mnula = [0]_{m\times n}$},
                                        com todos os elementos nulos
  \item[Matrix Identidade] \tabto{13mm} matriz com uns na diagonal
\end{description}
\[
  \onslide<+->{I_k}
  \onslide<+->{= \big[\delta_{ij}\big]_{k \times k}}
  \onslide<+->{= \begin{pmatrix}
    1      & 0      & \cdots & 0      \\
    0      & 1      & \cdots & 0      \\
    \vdots & \vdots & \ddots & \vdots \\
    0      & 0      & \cdots & 1
  \end{pmatrix}}
  \qquad\qquad
  \onslide<+->{\delta_{ij} = \begin{cases}
  1 & i = j \\
  0 & i \neq j
  \end{cases}}
\]
\end{frame}

%------------------------------------------------------------------------------%
\begin{frame}{Vetores e Escalares}
Uma matriz coluna é considerada um \emph{vetor}

\pause
\bigskip
\bigskip
Um número é chamado de \emph{escalar}
\end{frame}

%------------------------------------------------------------------------------%
%\section{Lista Mínima}
%
%%------------------------------------------------------------------------------%
%\begin{frame}{Lista Mínima}
%  \begin{enumerate}
%    \item
%  \end{enumerate}
%
%  \vfill
%  Atenção: A prova é baseada no livro, não nas apresentações
%\end{frame}

%------------------------------------------------------------------------------%
\end{document}
%------------------------------------------------------------------------------%
